\chapter{Endometriosis}
\index{Endometriosis}

\section*{Definition and Overview}
Endometriosis is a chronic condition in which tissue similar to the lining of the uterus grows outside the uterus, most commonly on the ovaries, fallopian tubes, and the tissue lining the pelvis. This tissue behaves like the uterine lining, thickening, breaking down, and bleeding with each menstrual cycle, but it has no way to leave the body. The result is inflammation, scarring, and often severe pelvic pain. Endometriosis affects around one in ten women and can affect fertility, but effective treatments are available for both pain and fertility.

\section*{Epidemiology}
Endometriosis affects approximately 10 per cent of women of reproductive age, making it one of the most common gynaecological conditions. It can begin in adolescence and persist until the menopause. Diagnosis is often delayed for many years, in part because period pain is normalised and because the condition requires surgery for definitive diagnosis. The condition has a significant impact on education, work, relationships, and quality of life.

\section*{Aetiology and Pathophysiology}
The cause of endometriosis is not fully understood, and several theories exist.
\begin{itemize}
    \item \textbf{Retrograde menstruation:} the most widely accepted theory, in which menstrual tissue flows backwards through the fallopian tubes and implants in the pelvis
    \item \textbf{Genetic factors:} the condition runs in families
    \item \textbf{Immune factors:} differences in the immune response may allow the tissue to implant and survive
    \item \textbf{Cellular transformation:} cells in the pelvis may transform into endometrial-like tissue
    \item \textbf{Effects:} the implants respond to hormones, causing inflammation, scarring, adhesions, and ovarian cysts called endometriomas
    \item \textbf{Pain mechanisms:} inflammation, nerve involvement, and adhesions cause pain that can continue beyond the period
    \item \textbf{Fertility:} scarring and inflammation can affect the tubes, ovaries, and egg quality
\end{itemize}

\section*{Clinical Features}
Symptoms vary widely, and some women have no symptoms at all.
\begin{itemize}
    \item Severe period pain that is not relieved by simple pain relief
    \item Chronic pelvic pain through the cycle
    \item Pain during or after intercourse
    \item Pain with bowel movements or urination
    \item Heavy or irregular periods
    \item Pain when opening the bowels during periods
    \item Difficulty becoming pregnant
    \item Fatigue, nausea, and bloating
    \item Symptoms may worsen over time and vary with the cycle
\end{itemize}

\section*{Differential Diagnosis}
Other conditions cause similar symptoms.
\begin{itemize}
    \item Primary dysmenorrhoea, period pain without an underlying condition
    \item Adenomyosis, endometriosis-like tissue within the uterine muscle
    \item Uterine fibroids and polyps
    \item Pelvic inflammatory disease
    \item Irritable bowel syndrome and other bowel conditions
    \item Ovarian cysts and ovarian cancer (rare)
    \item Musculoskeletal and urinary causes of pelvic pain
\end{itemize}

\section*{When to Seek Medical Care}
Women with severe period pain, pelvic pain, painful intercourse, or fertility concerns should be assessed by a GP and referred to a gynaecologist. Early diagnosis improves outcomes and quality of life.

\textbf{Contact your local emergency services for:}
\begin{itemize}
    \item Sudden severe pelvic or abdominal pain
    \item Heavy bleeding with faintness
    \item Fever with pelvic pain
    \item Shoulder tip pain with pelvic pain, which may indicate a ruptured cyst or ectopic pregnancy
\end{itemize}

\section*{Diagnosis and Investigations}
Diagnosis involves assessment and, for confirmation, surgery.
\begin{itemize}
    \item \textbf{History:} the pattern of pain, its relationship to the cycle, and its impact
    \item \textbf{Examination:} pelvic examination may reveal tenderness or nodules
    \item \textbf{Ultrasound:} detects ovarian endometriomas and deep disease, though normal ultrasound does not exclude endometriosis
    \item \textbf{Laparoscopy:} the definitive test, in which a camera is passed into the abdomen; endometriosis is visualised and can be treated at the same time
    \item \textbf{Trial of treatment:} response to hormonal treatment supports the diagnosis
\end{itemize}

\section*{Management}
Management is individualised, combining hormonal treatment, surgery, pain care, and lifestyle.
\begin{itemize}
    \item \textbf{Pain relief:} anti-inflammatory medicines such as ibuprofen
    \item \textbf{Hormonal contraception:} the combined pill, progestogen-only methods, and the hormonal IUD suppress the menstrual cycle and reduce pain
    \item \textbf{Other hormonal treatments:} GnRH analogues and danazol, which reduce oestrogen levels, with side effects to consider
    \item \textbf{Laparoscopic surgery:} removal or destruction of endometriosis tissue, which relieves pain and can improve fertility
    \item \textbf{Surgery for severe disease:} removal of the uterus and ovaries is a last resort for women who have completed their family
    \item \textbf{Fertility treatment:} assisted reproduction for women having difficulty conceiving
    \item \textbf{Multidisciplinary care:} physiotherapy, psychology, and diet, and support for the pain and its impact
\end{itemize}

\section*{Complications}
\begin{itemize}
    \item Chronic pelvic pain and its effects on work, study, and relationships
    \item Infertility
    \item Ovarian cysts, including large endometriomas
    \item Adhesions and scarring
    \item Depression and anxiety
    \item Effects of treatments, including surgical menopause
\end{itemize}

\section*{Prognosis and Outcomes}
Endometriosis is a chronic condition without a cure, but most women achieve good control with treatment. Pain responds well to hormonal suppression and surgery, although recurrence is common after surgery if hormonal suppression is not used. Fertility outcomes are good with assisted reproduction. Management is ongoing, and a team approach gives the best quality of life.

\section*{Prevention and Self-Care}
\begin{itemize}
    \item Seek early assessment for severe period pain
    \item Use hormonal contraception as part of the treatment plan if appropriate
    \item Exercise regularly, which reduces pain and inflammation
    \item Manage stress and seek psychological support
    \item Use heat, relaxation, and pacing during flare-ups
    \item Consider dietary changes with dietitian support
    \item Join support organisations for information and community
\end{itemize}

\section*{Sources and Further Reading}
The following reputable sources provide further information for healthcare students and patients seeking reliable, current advice about endometriosis.
\begin{itemize}
    \item Endometriosis Australia. \textit{Information and support.} https://www.endometriosisaustralia.org/
    \item Jean Hailes for Women's Health. \textit{Endometriosis.} https://www.jeanhailes.org.au/
    \item Healthdirect Australia. \textit{Endometriosis.} https://www.healthdirect.gov.au/endometriosis
    \item NHS. \textit{Endometriosis: overview and treatment.} https://www.nhs.uk/conditions/endometriosis/
    \item Mayo Clinic. \textit{Endometriosis: symptoms and causes.} https://www.mayoclinic.org/diseases-conditions/endometriosis/
    \item World Endometriosis Society. \textit{Patient resources.} https://endometriosis.org/
\end{itemize}
